% secciones/conclusiones.tex -- Sec 5 Conclusiones

Este estudio formuló y evaluó un modelo de caja gris para representar la generación de CO$_2$ durante la bioconversión con \textit{Hermetia illucens}, distinguiendo las fuentes larvaria y microbiana en cámaras estáticas. Se conservó la estructura del núcleo fisiológico y se verificó fórmula a fórmula contra su implementación publicada: E0 mantuvo los parámetros originales, mientras que C recalibró $a_{\max}$ y $\rho$ y añadió el módulo de prepupa.

En los datos propios, C redujo el RMSE censurado a 86.41~ppm/min en calibración y a 130.97~ppm/min en validación cruzada, frente a 180.53~ppm/min de E0 en esta última. La evidencia fuera de muestra es favorable, pero limitada: en los cuatro cierres exactos la reducción fue del 14.3~\%, la cobertura de las cotas saturadas fue 1/4 y el intervalo bootstrap de la mejora, $(-42.8,46.0)$~\%, incluye valores negativos. Los contrastes estructurales son compatibles con un control top-down ejercido mediante el pastoreo, pero no demuestran su identificabilidad biológica ni excluyen formulaciones alternativas.

La evaluación externa también es condicionada. En \textcite{ermolaevGreenhouseGasEmissions2019}, E0 reproduce el acumulado en transferencia pura; el nivel B de C reajusta cuatro parámetros con $n=p$, por lo que constituye una interpolación y presenta un compromiso entre el ajuste del gas y la biomasa final. En \textcite{schonUsingGasExchange2026}, C queda dentro del margen declarado en tres de cuatro dietas después de reajustar $f$ por dieta, mientras que la dieta más pobre permanece fuera. Estos resultados no sustentan una transportabilidad universal.

En consecuencia, el modelo es apropiado como representación de caja gris de este caso de estudio y mejora el desempeño global frente a la referencia, pero no ofrece una precisión predictiva uniforme ni una validez general demostrada. El trabajo futuro debe examinar la fracción asimilada del carbono pastoreado, incorporar réplicas biológicas independientes de gas, reducir la saturación instrumental y registrar la evolución dinámica de la humedad del lecho.
